\begin{appendices}

\section{Feature definitions, computation details and additional tables}\label{app:appendix}

\begin{table}[t]
\tabsize
\setlength{\tabcolsep}{3pt}
\caption{Observation suites for the three alternatives, three-cycle windows, common feature set, $\alpha=0.95$}\label{tab:suites3}
\begin{tabular}{@{}llllcccc@{}}
\toprule
Suite & Machine & Feature & MDE & DS & DS nested & Oracle DS & Oracle FA \\
\midrule
\multicolumn{8}{@{}l}{\emph{Inter-turn faults}} \\
drive-internal & PMSG & $D_2/D_1$ & 7.4 & 65 & 65 & 82 & 44 \\
drive-internal & SCIG & $I_d$\,2$f_e$ & 7.4 & 78 & 78 & 93 & 78 \\
drive-internal & WFSG & $I_q$\,2$f_e$ & 11.6$^{\mathrm{a}}$ & 91 & 87 & 96 & -- \\
3 CT & PMSG & $I_2/I_1$ & 7.4 & 57 & 57 & 78 & 11 \\
3 CT & SCIG & $I_2/I_1$ & 7.4 & 70 & 70 & 86 & 56 \\
3 CT & WFSG & $I_a$\,h3 & 11.6$^{\mathrm{a}}$ & 98 & 86 & 94 & -- \\
3 CT + 3 VT & PMSG & $V_2/V_1$ & 3.8 & 73 & 73 & 85 & 33 \\
3 CT + 3 VT & SCIG & $I_2/I_1$ & 7.4 & 70 & 70 & 89 & 56 \\
3 CT + 3 VT & WFSG & $V_2/V_1$ & 11.6$^{\mathrm{a}}$ & 98 & 98 & 94 & -- \\
3 CT + 3 VT + drive & PMSG & $V_2/V_1$ & 3.8 & 73 & 68 & 94 & 67 \\
3 CT + 3 VT + drive & SCIG & $I_d$\,2$f_e$ & 7.4 & 78 & 78 & 97 & 89 \\
3 CT + 3 VT + drive & WFSG & $V_2/V_1$ & 11.6$^{\mathrm{a}}$ & 98 & 98 & 98 & -- \\
excitation (field current) & WFSG & $I_f$\,2$f_e$ & 11.6$^{\mathrm{a}}$ & 86 & 86 & 74 & -- \\
\multicolumn{8}{@{}l}{\emph{Inter-winding faults}} \\
drive-internal & PMSG & $D_2/D_1$ & 14.3 & 91 & 91 & 98 & 44 \\
drive-internal & SCIG & $I_d$\,2$f_e$ & 2.3 & 100 & 100 & 100 & 78 \\
drive-internal & WFSG & $I_q$\,2$f_e$ & 2.3 & 100 & 100 & 100 & -- \\
3 CT & PMSG & $I_2/I_1$ & 9.7 & 83 & 83 & 99 & 11 \\
3 CT & SCIG & $I_2/I_1$ & 2.3 & 98 & 98 & 99 & 56 \\
3 CT & WFSG & $I_2/I_1$ & 2.3 & 100 & 100 & 100 & -- \\
3 CT + 3 VT & PMSG & $V_2/V_1$ & 9.7 & 92 & 92 & 99 & 33 \\
3 CT + 3 VT & SCIG & $I_2/I_1$ & 2.3 & 98 & 98 & 100 & 56 \\
3 CT + 3 VT & WFSG & $V_2/V_1$ & 2.3 & 100 & 100 & 100 & -- \\
3 CT + 3 VT + drive & PMSG & $V_2/V_1$ & 9.7 & 92 & 92 & 100 & 67 \\
3 CT + 3 VT + drive & SCIG & $I_d$\,2$f_e$ & 2.3 & 100 & 100 & 100 & 89 \\
3 CT + 3 VT + drive & WFSG & $I_q$\,2$f_e$ & 2.3 & 100 & 100 & 100 & -- \\
excitation (field current) & WFSG & $I_f$\,2$f_e$ & 2.3 & 100 & 100 & 100 & -- \\
\botrule
\end{tabular}
\footnotetext{Fixed feature: highest median SDR among the suite's screened features. MDE (\%) at the fault impedance matching the PMSG/SCIG protocol (WFSG: 198 trials, four tap pairs per class). DS: detectable share (\%) of the fixed feature; nested: fixed feature chosen on the other tap pairs; oracle: per-trial best feature, with its false-alarm rate on the healthy trials (FA, \%; not available for the WFSG). `--': largest tested extent not detectable.}
\footnotetext[\mathrm{a}]{The smallest inter-turn extent tested on the WFSG at that impedance is 11.6\,\%, so the value is an upper bound.}
\end{table}

\begin{table}[t]
\tabsize
\setlength{\tabcolsep}{3pt}
\caption{Sensitivity of MDE/DS to the window length (3, 5, 8 cycles at $\alpha=0.95$) and to the null quantile ($\alpha=0.90$ and $0.99$ at five cycles)}\label{tab:sens}
\begin{tabular}{@{}lllccccc@{}}
\toprule
Class & Feature & Machine & 3\,c & 5\,c & 8\,c & $\alpha$\,0.90 & $\alpha$\,0.99 \\
\midrule
IT & $V_2/V_1$ & PMSG & 3.8/73 & 3.8/76 & 3.8/73 & 3.8/80 & 3.8/69 \\
IT & $V_2/V_1$ & SCIG & 7.4/69 & 7.4/73 & 7.4/75 & 7.4/76 & 7.4/65 \\
IT & $I_2/I_1$ & PMSG & 7.4/57 & 7.4/56 & 7.4/58 & 7.4/65 & 7.4/52 \\
IT & $I_2/I_1$ & SCIG & 7.4/70 & 7.4/75 & 7.4/75 & 7.4/79 & 7.4/68 \\
IT & $I_d$\,2$f_e$ & PMSG & 12.1/29 & 12.1/30 & 12.1/28 & 12.1/33 & 12.1/24 \\
IT & $I_d$\,2$f_e$ & SCIG & 7.4/78 & 7.4/79 & 7.4/83 & 7.4/84 & 7.4/72 \\
IT & PI$_d$\,2$f_e$ & PMSG & 12.1/30 & 12.1/31 & 12.1/28 & 12.1/33 & 12.1/24 \\
IT & PI$_d$\,2$f_e$ & SCIG & 7.4/78 & 7.4/78 & 7.4/83 & 7.4/84 & 7.4/72 \\
IT & $V_q$\,2$f_e$ & PMSG & 7.4/55 & 7.4/57 & 7.4/60 & 7.4/65 & 7.7/50 \\
IT & $V_q$\,2$f_e$ & SCIG & 7.4/61 & 7.4/60 & 7.4/62 & 7.4/64 & 7.4/59 \\
IW & $V_2/V_1$ & PMSG & 9.7/92 & 2.3/94 & 2.3/94 & 2.3/94 & 9.7/91 \\
IW & $V_2/V_1$ & SCIG & 2.3/100 & 2.3/100 & 2.3/100 & 2.3/100 & 2.3/100 \\
IW & $I_2/I_1$ & PMSG & 9.7/83 & 9.7/81 & 9.7/82 & 3.2/87 & 9.7/79 \\
IW & $I_2/I_1$ & SCIG & 2.3/98 & 2.3/100 & 2.3/100 & 2.3/100 & 2.3/97 \\
IW & $I_d$\,2$f_e$ & PMSG & 21.7/71 & 21.7/72 & 21.7/69 & 21.7/74 & 21.7/69 \\
IW & $I_d$\,2$f_e$ & SCIG & 2.3/100 & 2.3/100 & 2.3/100 & 2.3/100 & 2.3/100 \\
IW & PI$_d$\,2$f_e$ & PMSG & 21.7/71 & 21.7/72 & 21.7/70 & 21.7/74 & 21.7/69 \\
IW & PI$_d$\,2$f_e$ & SCIG & 2.3/100 & 2.3/100 & 2.3/100 & 2.3/100 & 2.3/100 \\
IW & $V_q$\,2$f_e$ & PMSG & 9.7/89 & 9.7/90 & 9.7/91 & 9.7/91 & 9.7/86 \\
IW & $V_q$\,2$f_e$ & SCIG & 9.7/88 & 9.7/92 & 9.7/88 & 9.7/95 & 9.7/84 \\
\botrule
\end{tabular}
\footnotetext{Cells: MDE (\%; limit-of-detection convention; `--': largest extent not detectable) / DS (\%). IT: inter-turn; IW: inter-winding.}
\footnotetext[\mathrm{a}]{The feature fails the reliability screen at that setting.}
\end{table}

\begin{table}[t]
\tabsize
\setlength{\tabcolsep}{2.5pt}
\caption{Within-trial versus commissioning-baseline null, five-cycle windows, $\alpha=0.95$}\label{tab:between}
\begin{tabular}{@{}lllcccccc@{}}
\toprule
\multirow{2}{*}{Class} & \multirow{2}{*}{Feature} & \multirow{2}{*}{M.} & \multicolumn{2}{c}{$q_{95}$} & \multicolumn{2}{c}{Median SDR} & \multicolumn{2}{c}{DS} \\
\cmidrule(lr){4-5}\cmidrule(lr){6-7}\cmidrule(l){8-9}
 & & & within & baseline & within & baseline & within & baseline \\
\midrule
IT & $V_2/V_1$ & P & 0.086 & 0.315 & 2.99 & 1.05 & 76 & 50 \\
IT & $I_2/I_1$ & P & 0.162 & 0.268 & 1.47 & 0.978 & 56 & 48 \\
IT & $I_d$\,2$f_e$ & P & 0.563 & 0.876 & 0.516 & 0.441 & 30 & 20 \\
IT & PI$_d$\,2$f_e$ & P & 0.557 & 0.870 & 0.518 & 0.441 & 31 & 20 \\
IT & $V_q$\,2$f_e$ & P & 0.165 & 0.277 & 1.35 & 0.872 & 57 & 44 \\
IT & $D_2/D_1$ & P & 0.090 & 0.211 & 2.78 & 1.11 & 64 & 52 \\
IW & $V_2/V_1$ & P & 0.086 & 0.315 & 61.9 & 12.5 & 94 & 83 \\
IW & $I_2/I_1$ & P & 0.162 & 0.268 & 11.3 & 6.34 & 81 & 75 \\
IW & $I_d$\,2$f_e$ & P & 0.563 & 0.876 & 4.14 & 2.4 & 72 & 67 \\
IW & PI$_d$\,2$f_e$ & P & 0.557 & 0.870 & 4.21 & 2.4 & 72 & 67 \\
IW & $V_q$\,2$f_e$ & P & 0.165 & 0.277 & 22.4 & 10.7 & 90 & 81 \\
IW & $D_2/D_1$ & P & 0.090 & 0.211 & 32.6 & 11.9 & 90 & 85 \\
IT & $V_2/V_1$ & S & 0.033 & 0.111 & 4.17 & 1.45 & 73 & 59 \\
IT & $I_2/I_1$ & S & 0.034 & 0.130 & 4.52 & 1.15 & 75 & 54 \\
IT & $I_d$\,2$f_e$ & S & 0.016 & 0.078 & 6.68 & 1.29 & 79 & 53 \\
IT & PI$_d$\,2$f_e$ & S & 0.016 & 0.078 & 6.69 & 1.29 & 78 & 53 \\
IT & $V_q$\,2$f_e$ & S & 0.060 & 0.131 & 2.58 & 1.54 & 60 & 58 \\
IT & $D_2/D_1$ & S & 0.010 & 0.068 & 11.1 & 1.57 & 81 & 52 \\
IW & $V_2/V_1$ & S & 0.033 & 0.111 & 48.2 & 16.1 & 100 & 94 \\
IW & $I_2/I_1$ & S & 0.034 & 0.130 & 55.6 & 12.9 & 100 & 86 \\
IW & $I_d$\,2$f_e$ & S & 0.016 & 0.078 & 74 & 16 & 100 & 88 \\
IW & PI$_d$\,2$f_e$ & S & 0.016 & 0.078 & 74 & 16 & 100 & 88 \\
IW & $V_q$\,2$f_e$ & S & 0.060 & 0.131 & 24.8 & 11.2 & 92 & 86 \\
IW & $D_2/D_1$ & S & 0.010 & 0.068 & 84.9 & 12.2 & 100 & 82 \\
\botrule
\end{tabular}
\footnotetext{Within: reference immediately before the fault; baseline: the healthy trial at the same operating point, recorded at another time. DS: detectable share (\%). IT: inter-turn; IW: inter-winding; P: PMSG; S: SCIG.}
\end{table}

\begin{table}[t]
\tabsize
\caption{Transfer between the three alternatives with a self-referenced GBDT detector, three-cycle windows, common features}\label{tab:transfer3}
\begin{tabular}{@{}llcccc@{}}
\toprule
Trained on & Tested on & AUC all & AUC o.o.r. & TPR o.o.r. & FA healthy (\%) \\
\midrule
PMSG & PMSG & 0.86 & 0.85 & 0.62 & 3 \\
SCIG & SCIG & 0.92 & 0.91 & 0.76 & 1 \\
WFSG & WFSG & 0.97 & 0.97 & 0.93 & -- \\
PMSG & SCIG & 0.88 & 0.88 & 0.60 & 0 \\
PMSG & WFSG & 0.97 & 0.98 & 0.89 & -- \\
SCIG & PMSG & 0.81 & 0.80 & 0.35 & 0 \\
SCIG & WFSG & 0.95 & 0.96 & 0.86 & -- \\
WFSG & PMSG & 0.86 & 0.84 & 0.52 & 1 \\
WFSG & SCIG & 0.91 & 0.89 & 0.61 & 0 \\
SCIG+WFSG & PMSG & 0.84 & 0.83 & 0.41 & 0 \\
PMSG+WFSG & SCIG & 0.89 & 0.88 & 0.61 & 0 \\
PMSG+SCIG & WFSG & 0.97 & 0.97 & 0.88 & -- \\
\botrule
\end{tabular}
\footnotetext{Features $|z|$-referenced to the first half of each trial's own reference interval. AUC on all windows and on the out-of-reference (o.o.r.) windows; TPR: true-positive rate at 1\,\% false-positive rate (FPR) on the o.o.r. windows; FA: share of the healthy trials' fault-time windows above the threshold giving 1\,\% FPR on all negatives. The WFSG has no healthy trials, so its negatives are reference and recovery windows only and its rows are optimistic.}
\end{table}


This appendix defines the features (Section~\ref{app:features}), lists the computation details (Section~\ref{app:computation}) and gives the observation-suite results for the three alternatives (Table~\ref{tab:suites3}), the sensitivity of the design quantities to the window length and the null quantile (Table~\ref{tab:sens}), the commissioning-baseline null (Table~\ref{tab:between}) and the transfer results between the three alternatives (Table~\ref{tab:transfer3}).

\subsection{Feature definitions}\label{app:features}

Let $x_a, x_b, x_c$ be the sampled phase quantities of a window of $N$ samples at sampling frequency $f_s$, $w$ the Hann window and $f_e$ the electrical frequency of the window, estimated from the slope of the unwrapped phase of the current space vector $\tfrac{2}{3}(i_a + \mathbf{a}\, i_b + \mathbf{a}^2 i_c)$, $\mathbf{a}=e^{j2\pi/3}$, whose sign gives the phase rotation (outside 0.5--1.5 times the nominal value the nominal value is used). The complex fundamental amplitude of a signal $x$ at frequency $f$ is $X(f) = \frac{2}{\sum_n w_n}\sum_{n=0}^{N-1} x_n\, w_n\, e^{-j2\pi f n/f_s}$. Sequence magnitudes are $|X_+| = |X_a + \mathbf{a}X_b + \mathbf{a}^2X_c|/3$ and $|X_-| = |X_a + \mathbf{a}^2X_b + \mathbf{a}X_c|/3$ for an $abc$ rotation (interchanged for $acb$), evaluated at $f=f_e$ for currents ($I_1, I_2$), terminal voltages ($V_1, V_2$) and mean-removed duty cycles ($D_1, D_2$); $|X(f)|$ is a peak amplitude and is divided by $\sqrt2$ wherever it is compared with an RMS value. The RMS unbalance is $\max_k |I_{\mathrm{rms},k} - \bar I_{\mathrm{rms}}| / \bar I_{\mathrm{rms}}$ over the three phases; the harmonic ratios are $|I_a(kf_e)|/|I_a(f_e)|$ for $k=3,5,7$ and the THD is $\sqrt{\sum_{k=2}^{20}|I_a(kf_e)|^2}/|I_a(f_e)|$. For the synchronous-frame currents $i_d, i_q$ and the torques the features are the standard deviation of the mean-removed signal and its amplitudes $|Y(f_e)|$ and $|Y(2f_e)|$; for the PI outputs the standard deviation and $|Y(2f_e)|$; for the measured and commanded synchronous-frame voltages $|Y(2f_e)|$; for the field current the standard deviation and $|Y(2f_e)|$; speed and DC-link features are standard deviations. Windows are $n_c$ electrical cycles long with a hop of one cycle, and a feature is undefined for a trial whose faulted interval holds no complete window. On the WFSG bench the measured converter-side phase voltages serve as terminal voltages, the synchronous-frame voltages are obtained from them with the encoder angle, the commanded voltages are the controller's, and the electrical torque is computed from power and speed.

\subsection{Computation details}\label{app:computation}

\emph{Intervals.} PMSG and SCIG: reference interval 0.30--0.98\,s after the start of each record, split at 0.62\,s into $R_1$ (0.32\,s) and $R_2$ (0.36\,s); contactor commanded at 1.000\,s; faulted interval from 30\,ms after the command to 10\,ms before its release; the relay picks up 85--90\,ms after the command, so the faulted interval holds about 55\,ms of healthy time (at most 15\,\% of its length), which dilutes the fault-interval mean and makes the PMSG/SCIG SDRs slightly conservative; recovery windows start 0.30\,s after the release. WFSG: fault interval located from the measured fault current as the first contiguous interval in which its 10\,ms RMS exceeds $\max(0.3\,\mathrm{A}, 6\sigma)$, $\sigma$ being the standard deviation of that RMS over 0.05--0.40\,s, ending at the first fall below the threshold after 50\,ms; reference interval from 0.05\,s after the start of the record to 10\,ms before the onset, split at 0.30\,s; recovery windows start 0.15\,s after the end of the fault interval. At 23.9--63.7\,Hz a WFSG faulted interval holds between zero and eleven three-cycle windows; four trials hold none and are dropped, and 81 of 438 hold fewer than three (all 50 trials at 1671\,rpm, 26 of 206 at 1800\,rpm, 4 at 716 and 1 at 955\,rpm).

\emph{Null and screen.} The null is pooled over all 225 records of each of the PMSG and SCIG and over the 438 inter-turn and inter-winding records of the WFSG. The screen flags a feature when its count of false alarms on the nine healthy trials is incompatible with the nominal rate $1-\alpha$ at the 5\,\% level (one-sided binomial test: three or more at $\alpha=0.95$, two or more at $\alpha=0.99$, four or more at $\alpha=0.90$) and excludes features whose null quantile exceeds one. At $\alpha=0.95$ it removes, at three cycles, $D_2/D_1$ on the SCIG; at five cycles $I_d$\,1$f_e$ and speed std on the PMSG and $I_a$\,h7 and $D_2/D_1$ on the SCIG; at eight cycles PI$_d$\,std and speed std on the PMSG and $I_q$\,2$f_e$, $D_2/D_1$ and $V_{dc}$\,std on the SCIG.

\emph{Bootstrap.} 2000 resamples, seed 1. Median SDR and detectable share of a class: cluster bootstrap over the twelve tap pairs, the same resampled pairs for both machines; Wilson 95\,\% intervals for shares. Minimum detectable extent: operating-point trials resampled within each extent level and the limit-of-detection rule applied to each resample. The null quantile is held at its full-sample value. Decision table: 1000 resamples per suite, alternative, class and $\alpha$ (seed 7); a suite \emph{meets} $e_{\mathrm{req}}$ when the resample probability of $\mathrm{MDE}\le e_{\mathrm{req}}$ is at least 0.95 and \emph{fails} when it is at most 0.05.

\emph{Transfer.} Models: scikit-learn histogram gradient boosting (maximum depth 3, 200 iterations, learning rate 0.08, random state 0) and logistic regression with standardisation ($C=1$, 3000 iterations); five-fold GroupKFold with one group per tap pair and machine, healthy trials grouped by speed. The self-referencing rule uses the windows starting before 0.62\,s (PMSG/SCIG) or 0.30\,s (WFSG) as the normalising set. Metrics: area under the receiver operating characteristic curve (AUC) and true-positive rate at 1\,\% false-positive rate, on all windows and on the windows outside the normalising set, and the false-alarm rate of the all-negatives 1\,\% threshold on the fault-time windows of the healthy trials.

\end{appendices}
